%%% Основные сведения %%%
\newcommand{\thesisAuthor}             % Диссертация, ФИО автора
{%
    \texorpdfstring{% \texorpdfstring takes two arguments and uses the first for (La)TeX and the second for pdf
        Бабкин Пётр % так будет отображаться на титульном листе или в тексте, где будет использоваться переменная
    }{%
        Бабкин Пётр % эта запись для свойств pdf-файла. В таком виде, если pdf будет обработан программами для сбора библиографических сведений, будет правильно представлена фамилия.
    }%
}

\newcommand{\thesisTitle}              % Диссертация, название
{\texorpdfstring{\MakeUppercase{Поиск архитектуры модели смеси экспертов с учётом кластерной структуры данных}}{Поиск архитектуры модели смеси экспертов с учётом кластерной структуры данных}}
\newcommand{\thesisSpecialtyNumber}    % Диссертация, специальность, номер
{\texorpdfstring{09.04.01}{09.04.01}}
\newcommand{\thesisSpecialtyTitle}     % Диссертация, специальность, название
{\texorpdfstring{Информатика и вычислительная техника}{Информатика и вычислительная техника}}
\newcommand{\thesisDegree}             % Диссертация, научная степень
{Выпускная квалификационная работа магистра}
\newcommand{\thesisCity}               % Диссертация, город защиты
{Москва}
\newcommand{\thesisYear}               % Диссертация, год защиты
{2026}
\newcommand{\thesisOrganization}       % Диссертация, организация
{Федеральное государственное автономное образовательное учреждение
 высшего образования
 «Московский физико-технический институт
(национальный исследовательский университет)»}

\newcommand{\supervisorFio}            % Научный руководитель, ФИО
{Бахтеев О.Ю.}

\newcommand{\keywords}%                 % Ключевые слова для метаданных PDF диссертации и автореферата
{смесь экспертов · поиск нейросетевой архитектуры · суррогатная модель · EM-алгоритм · кластеризация данных}